\documentclass[18px, A4]{article}

\usepackage{tikz}
\usepackage{amsmath}

\begin{document}
\title{Numerische Simulation lichtartiger Geodähten in der Schwarzschild-Metrik}
\author{Beat Grüner}
\pagenumbering{gobble}
\maketitle
\newpage
\pagenumbering{arabic}

\tableofcontents
\newpage



\section{Harmonischer Oszillator als motivierendes Beispiel}
Ein Standardbeispiel für einfache Differentialgleichungen ist der harmonische Oszilator.
Der Harmonische Oszillator kann durch Pendel oer elektronische Schwingkreise motiviert werden.
Der ungedämpfte harmonische Oszillator ist die Lösung der DGL
\[
    \frac{d^{2}x}{dt^{2}}\ +\ \omega^{2} x\ =\ 0
\]

\subsection{Umschreiben in ein System erster Ordnung}
Diese DGL kann in ein System erster Ordnung umgeschrieben werden. Dafür wird der
Zustand des Systems als Vektor
\[
    \vec{y}\ =\ \begin{bmatrix}
	x  \\
	x' \\
    \end{bmatrix}
\]
geschrieben. \newline
Wenn wir nun die DGL wie folgt umstellen:
\[
    \frac{d^{2}x}{dt^{2}}\ =\ -\omega^{2} x
\]
können wir nun die vektorwertige Funktion
\[
    \vec{y}\ '\ =\ f(\vec{y}, t)\ =\ \begin{bmatrix}
	x' \\
	-\omega^{2} x
    \end{bmatrix}
\]
definieren, da $x''\ =\ -\omega x$.

\subsection{Exakte Lösung des Problems}
Die Form der allgemeine Lösung der DGL ist:
\[
    x(t)\ =\ \sin \omega t
\]
wobei $\omega$ die Kreisfrequenz der Oszillation ist. 

\newpage


\section{Numerische Verfahren angewandt auf den harmonischen Oszillator}
Auf dieses System erster Ordnung können nun verschiedene numerische Verfahren angewandt werden

\subsection{Explizites Euler-verfahren}
Das explizite Euler-verfahren

\subsection{runge-kutta-verfahren}

\subsection{implizites euler-verfahren}

\subsection{implizite mittelpunktregel (symplektisch)}

\newpage


\section{Die Geodähtengleichung auf \(S^{2}\)}

\subsection{Die Geodähtengleichung auf \(S^{2}\) als System erster Ordnung}

\subsection{Explizites Euler-Verfahren}

\subsection{Runge-Kutta-Verfahren}

\subsection{Implizites Euler-Verfahren}

\subsection{Implizite Mittelpunktregel (Symplektisch)}

\newpage


\section{Die Schwarzschild-Metrik}

\subsection{Die Schwarzschild-Metrik als System erster Ordnung}

\subsection{Vereinfachung des Systems}

\subsection{Explizites Euler-Verfahren}

\subsection{Runge-Kutta-Verfahren}

\subsection{Implizites Euler-Verfahren}
\subsubsection{Verfahren zum Lösen von Gleichungssystemen}

\subsection{Implizite Mittelpunktregel (Symplektisch)}



\end{document}
